\section{Detecção e Tratamento de \glspl{vm} em Idle na RISC-V}

Em plataformas RISC-V, existe uma instrução específica - \gls{wfi}, utilizada para colocar um núcleo da CPU em estado de \textit{idle}. Ao executá-la, o processador suspende a execução normal até a
ocorrência de uma interrupção, reduzindo o consumo de energia enquanto não há tarefas a serem processadas.

Em ambientes virtualizados, é necessário que o hipervisor intercepte essa instrução para gerenciar corretamente o comportamento dos sistemas convidados.

Para isso, a arquitetura disponibiliza um mecanismo de controle de acesso a essa instrução por meio do bit 21 (\textit{Timeout Wait}) do \gls{csr} \textit{mstatus}.

Esse bit de configuração permite tratar a instrução \gls{wfi} como uma \textit{illegal instruction}.
Ou seja, quando esse bit está ativado, a execução de uma \gls{wfi} não resulta diretamente
no estado de espera do processador; em vez disso, uma exceção é gerada.
Esta exceção é então detectada pelo hipervisor, que - em seguida - faz com que essa \gls{vm} se afaste do controle da CPU, permitindo que outra ocupe seu lugar. No fim, somente se nenhuma outra
\textit{Thread} estiver disponível, a CPU entrará verdadeiramente em \textit{idle}.
